% =====================================================================
%  THÈME 2 — Conversions en quantité de matière
%  Niveau MOYEN — Exercices
% =====================================================================
\documentclass[11pt,a4paper]{article}
\usepackage[utf8]{inputenc}
\usepackage[T1]{fontenc}
\usepackage{lmodern}
\usepackage[french]{babel}
\usepackage{amsmath,amssymb,mathtools}
\usepackage{icomma}
\usepackage[version=4]{mhchem}
\usepackage{siunitx}
\sisetup{output-decimal-marker={,},per-mode=symbol}
\DeclareSIUnit\Molar{\mole\per\litre}
\DeclareSIUnit\atm{atm}
\usepackage{xcolor}
\usepackage{tikz}
\usepackage[most]{tcolorbox}
\usepackage{enumitem}
\usepackage{array}
\usepackage{fancyhdr}
\usepackage[a4paper,margin=2.1cm,headheight=26pt,headsep=8pt]{geometry}

\definecolor{primary}{RGB}{142,93,168}
\definecolor{primarydark}{RGB}{82,45,108}
\definecolor{moycol}{RGB}{198,120,9}
\newcommand{\gmol}{\si{\gram\per\mole}}
\newcommand{\Vm}{V_{\mathrm m}}

\pagestyle{fancy}\fancyhf{}
\fancyhead[L]{\small\textcolor{primarydark}{\textbf{Fiche 5} — Thème 2 : Conversions}}
\fancyhead[R]{\small\textcolor{moycol}{\textsc{Moyen}}}
\fancyfoot[C]{\small\thepage}
\renewcommand{\headrulewidth}{0.8pt}
\renewcommand{\headrule}{\hbox to\headwidth{\color{primary}\leaders\hrule height \headrulewidth\hfill}}

\newtcolorbox[auto counter]{exo}[1][]{enhanced,breakable,boxrule=1pt,arc=3pt,
  left=3mm,right=3mm,top=2.4mm,bottom=2.4mm,colback=moycol!6,colframe=moycol,
  fonttitle=\bfseries,coltitle=white,
  title={Exercice~\thetcbcounter\ \ifx\relax#1\relax\relax\else\textmd{--- \itshape #1}\fi}}

\setlength{\parindent}{0pt}
\setlength{\parskip}{3pt}

\begin{document}

\begin{center}
{\LARGE\bfseries\color{primarydark}Thème 2 — Conversions}\\[1pt]
{\large\color{moycol}\textbf{Niveau Moyen}}\\[2pt]
{\itshape Objectif : enchaîner deux conversions (masse, volume, densité, entités).}\\
{\color{primary}\rule{0.6\linewidth}{1pt}}
\end{center}

\textbf{Données :} $N_A=\SI{6,02e23}{}$ ; $\Vm=\SI{22,4}{\litre\per\mole}$ (TPN).\quad
Masses molaires (\gmol) fournies dans chaque énoncé.

\begin{exo}[du gramme au volume gazeux]
On dispose de \SI{16}{\gram} de dioxygène \ce{O2} ($M=\SI{32}{\gram\per\mole}$).
Quel volume ce gaz occupe-t-il aux conditions TPN ?
\emph{(Chaîne : masse $\xrightarrow{\div M}$ mole $\xrightarrow{\times \Vm}$ volume.)}
\end{exo}

\begin{exo}[la densité d'un liquide]
Le dichlorométhane \ce{CH2Cl2} ($M=\SI{85}{\gram\per\mole}$) a une densité $d=1{,}70$
(c.-à-d.\ $\SI{1,70}{\gram\per\milli\litre}$). On en prélève \SI{100}{\milli\litre}.
\textbf{a)} Quelle est la masse prélevée ?\quad
\textbf{b)} À combien de moles cela correspond-il ?
\end{exo}

\begin{exo}[peser un nombre voulu de molécules]
Le mannitol ($M=\SI{182}{\gram\per\mole}$) est un sucre utilisé en perfusion. On souhaite disposer
d'exactement $\SI{1,20e21}{}$ molécules de mannitol.
\textbf{a)} À quelle quantité de matière (en mol) cela correspond-il ?\quad
\textbf{b)} Quelle masse (en mg) faut-il alors peser ?
\end{exo}

\begin{exo}[du comprimé à la concentration]
Un comprimé apporte \SI{755}{\milli\gram} de paracétamol \ce{C8H9NO2}
($M=\SI{151}{\gram\per\mole}$), que l'on dissout dans \SI{200}{\milli\litre} d'eau.
\textbf{a)} Combien de moles de paracétamol le comprimé contient-il ?\quad
\textbf{b)} Quelle est la concentration molaire de la solution obtenue ?
\end{exo}

\begin{exo}[remonter à la masse molaire d'un gaz]
Un échantillon de \SI{5,0}{\gram} d'un gaz inconnu occupe \SI{5,6}{\litre} aux conditions TPN.
\textbf{a)} Quelle est sa quantité de matière ?\quad
\textbf{b)} En déduire sa masse molaire $M$. \emph{(Chaîne : volume $\to$ mole, puis $M=m/n$.)}
\end{exo}

\end{document}
